% !TeX root = closed-systems-from-comparison-completeness.tex
% ============================================================
\chapter*{Orientation to Part IV}
\phantomsection
\label{part:intro-dynamics-spacetime}
% ============================================================
\begin{chapterorientation}
\orientationitem{Settled}{Part~III identifies the intrinsic transport obstruction and the two admissible enrichment loci.}
\orientationitem{Task}{Part~IV introduces internal observer reduction and interleaves it with refinement inside the closed system.}
\orientationitem{Handoff}{The output is the stitched arena \(ST\), which provides the comparison-side setting for curvature realization.}
\end{chapterorientation}

Part~IV does not introduce an external clock or an external coarse
graining.  It asks how irreversibility and spacetime-like organization
can appear inside a quotient theory whose underlying representative
dynamics remains controlled by the transport stack.

\Cref{ch:flow} isolates the conditional observer-level mechanism by
which irreversible descent can arise without changing the quotient
semantics.  \Cref{ch:stitching} then combines observer reduction with
the refinement tower to form the intrinsic spacetime arena \(ST\).  The
part supplies the stage on which loop defect can be read as curvature.
